\documentclass[10pt,a4paper]{amsart}
\usepackage[margin=19mm]{geometry}
\usepackage{amsmath,amssymb,mathtools}
\usepackage[scr=rsfs,cal=euler]{mathalfa}
\usepackage{xfrac}
\usepackage[unicode,hidelinks,bookmarks=false]{hyperref}
\newcommand{\ie}{i.e.\ }
\newcommand{\cf}{cf.\ }
\newcommand{\resp}{respectively\ }
\newcommand{\Subsection}{\subsection}
\providecommand{\nobreakdash}{\nobreak\hbox{-}}
\setlength{\emergencystretch}{2em}
\multlinegap0pt
\usepackage{bm}
\usepackage{bbm}
\usepackage{dsfont}
\usepackage{xspace}
\usepackage{cancel}
\usepackage{mathtools}
\usepackage{amsthm}
\newtheorem{theorem}{Theorem}
\newtheorem{lemma}[theorem]{Lemma}
\title[On convex bodies with rotationally symmetric planar projecti]{On convex bodies with rotationally symmetric planar projections}
\author{Sergii Myroshnychenko, Dmitry Ryabogin,, Christos Saroglou}
\date{Source: arXiv:2609.27192v1 (CC BY 4.0)}
\begin{document}
\maketitle

\noindent\emph{Source and licence.} On convex bodies with rotationally symmetric planar projections. Version arXiv:2609.27192v1 (CC BY 4.0), distributed under the Creative Commons Attribution 4.0 International licence, as recorded in the arXiv licence metadata for this exact version and shown on the abstract page https://arxiv.org/abs/2609.27192v1. Full rights notice: licenses/arxiv-2609.27192-CC-BY-4.0.txt.\par\smallskip

\noindent\textit{Editorial scope.} Counted unit: the Lemma whose source label is \texttt{l1O177} in the article (source0.tex lines 1225--1236, source label \texttt{l1O177}), delivered together with the article's own complete original proof of it (source0.tex lines 1238--1345, one \texttt{proof} environment of 108 lines). No mathematics is written, repaired or re-derived: statement and proof are the article's own text line for line. Required context retained uncounted: none. Every definition, display and label of those lines is retained unchanged. Omitted from this package and not redistributed: 1--1224, 1237--1237, 1346--1355. Nothing inside a retained excerpt is omitted or altered: every retained body is delivered exactly as the article's lines are written, so no omission receipt of any kind accompanies this package. Citations: the retained excerpts carry no citation command at all, so no bibliography entry of the article is transcribed and no reference is redistributed. Reference adaptation: none was needed and none was made; every reference inside the delivered unit resolves inside it, so no unresolved reference or citation is produced. Printed slips kept literal: no printed word, symbol, hypothesis or number of the counted statement or of its proof was altered. Layout and class: the article's own class is \texttt{article} with packages including amsfonts, latexsym, amsmath, amssymb, color, amsthm, fullpage, fontenc, inputenc, enumerate; the wrapper is the lane's pinned \texttt{amsart} editorial wrapper at 10pt on A4, which restates the theorem-like environments the retained text needs and adds the article's own macro definitions that the retained text needs (none), with extra wrapper packages (bm, bbm, dsfont, xspace, cancel, mathtools). The delivered numbering is the unit's own. Assets: no retained span contains a figure environment, a graphics-inclusion command, a PDF-page inclusion, a table or a TikZ picture; no image file is copied into this package, the package declares no asset dependency and no image file is redistributed with it. Licence: version arXiv:2609.27192v1 of this article is distributed under the Creative Commons Attribution 4.0 International licence, as recorded in the arXiv licence metadata for this exact version and shown on the abstract page https://arxiv.org/abs/2609.27192v1. A full rights notice is delivered at licenses/arxiv-2609.27192-CC-BY-4.0.txt. Characters: every retained excerpt is pure ASCII, so the delivered engine, XeLaTeX with its default Latin Modern text font, reports no missing character.
\begin{lemma}\label{l1O177}
	Let \(Y\in\mathcal H_\ell(S^2)\), \(Y\neq0\), and let \(r\in\mathbb Z\)
	satisfy
	\[
	|r|\leq \ell,
	\qquad
	r\equiv \ell\pmod 2.
	\]
	Then there exists a great circle \(\mathcal C\subset S^2\), with a suitable
	orientation and angular parametrization, such that the \(r\)-th Fourier
	coefficient of \(Y|_{\mathcal C}\) is nonzero.
\end{lemma}

\begin{proof}
	Let
	\[
	\mathcal C_0=\{z=0\}\cap S^2
	\]
	be the standard equator, parametrized by
	\[
	u(\varphi)=(\cos\varphi,\sin\varphi,0),
	\qquad
	0\leq\varphi<2\pi.
	\]
	Define the linear functional
	\[
	T_r:\mathcal H_\ell(S^2)\longrightarrow\mathbb C
	\]
	by
	\[
	T_r(Z)
	=
	\frac{1}{2\pi}
	\int_0^{2\pi}
	Z(u(\varphi))e^{-ir\varphi}\,d\varphi.
	\]
	Thus \(T_r(Z)\) is the \(r\)-th Fourier coefficient of
	\(Z|_{\mathcal C_0}\).
	
	We first show that \(T_r\not\equiv0\). Let \(Y_\ell^r\) be a standard
	spherical harmonic of degree \(\ell\) and order \(r\). For \(r\geq0\), it
	has the form
	\[
	Y_\ell^r(\theta,\varphi)
	=
	c_{\ell,r}P_\ell^r(\cos\theta)e^{ir\varphi}.
	\]
	Restricting to the equator \(\theta=\pi/2\), we obtain
	\[
	Y_\ell^r\Bigl(\frac{\pi}{2},\varphi\Bigr)
	=
	c_{\ell,r}P_\ell^r(0)e^{ir\varphi}.
	\]
	It is well known that
	\[
	P_\ell^r(0)\neq0
	\qquad\Longleftrightarrow\qquad
	\ell-r\ \text{is even}.
	\]
	Since \(r\equiv\ell\pmod2\), the right-hand side is nonzero, and hence
	\[
	T_r(Y_\ell^r)\neq0.
	\]
	
	For \(r<0\), we use the standard relation
	\[
	Y_\ell^r
	=
	c\,\overline{Y_\ell^{-r}}
	\]
	with \(c\neq0\). Its restriction to the equator is therefore a nonzero
	multiple of \(e^{ir\varphi}\), because
	\[
	\ell-(-r)=\ell-|r|
	\]
	is even. Thus
	\[
	T_r(Y_\ell^r)\neq0
	\]
	also in this case. Consequently,
	\[
	T_r\not\equiv0
	\]
	for every admissible \(r\).
	
	For \(R\in SO(3)\), let
	\[
	(U_RZ)(u)=Z(R^{-1}u).
	\]
	Suppose, contrary to the assertion, that the \(r\)-th Fourier coefficient of
	\(Y\) vanishes on every rotated equator, with the parametrization induced by
	the rotation. Equivalently,
	\[
	T_r(U_RY)=0
	\qquad
	\text{for every }R\in SO(3).
	\]
	Therefore \(T_r\) vanishes on the linear span of the orbit
	\[
	\{U_RY:R\in SO(3)\}.
	\]
	
	The span of this orbit is a nonzero \(SO(3)\)-invariant subspace of
	\(\mathcal H_\ell(S^2)\). Since \(\mathcal H_\ell(S^2)\) is irreducible, this
	span is the whole space:
	\[
	\operatorname{span}\{U_RY:R\in SO(3)\}
	=
	\mathcal H_\ell(S^2).
	\]
	It follows that \(T_r\equiv0\) on \(\mathcal H_\ell(S^2)\), contradicting
	the fact proved above that \(T_r\not\equiv0\).
	
	Hence, for some \(R\in SO(3)\),
	\[
	T_r(U_RY)\neq0.
	\]
	Equivalently, the restriction of \(Y\) to the great circle
	\(R^{-1}\mathcal C_0\), with the parametrization induced by \(R\), has a
	nonzero \(r\)-th Fourier coefficient.
\end{proof}

\end{document}
